\begin{theorem}[Main Theorem] \label{thm:main} For any fixed
prime $q \geq 5$ and constant $\zeta > 0$,
it is \cclass{NP}-hard to distinguish whether a 2P1R Game
with $R = q^6$ answers has value at least $1-\zeta$ or at most ${4}/{q}$.
\end{theorem}

\subsection{Overview of proofs and techniques}

We prove \expref{Theorem}{thm:main} by constructing a PCP (for an \cclass{NP}-complete language)
that makes two queries to a proof: one query reads a symbol over an alphabet of
size $q$ and the other reads a symbol
over an alphabet of size $q^6$. The PCP has completeness $1-\zeta$
and soundness at most ${4}/{q}$.
As is standard in the PCP literature, the PCP is obtained by composing the so-called
Outer PCP and Inner PCP.

\paragraph{Inner PCP}
The Inner PCP is essentially a probabilistic
testing procedure that tests whether a given function $f: \F_q^m \mapsto \F_q$
satisfies a desired property. Three types of tests have generally been used depending
on the desired setting of parameters, \eg, the number of queries, type of the
acceptance predicate,
completeness and soundness, size of the proof, etc.
\begin{itemize}
\item The low degree test~\cite{FGLSS, AS, AroraSudan, RazSafra} that tests whether $f$ is a polynomial of low degree.
\item The linearity test that tests whether $f$ is linear (= Hadamard codeword) \cite{ALMSS, Kh-maxclique}.
\item The dictatorship test that tests whether $f$ is a dictatorship, \ie, a function
of the form $f(x)=x_i$ for some coordinate $1 \leq i \leq m$, \eg,~\cite{BGS, Has96, Has97}.
\end{itemize}
If $f$ satisfies the desired property (\ie, being linear, low degree, or dictatorship),
then the test passes with probability
(close to) $1$ and conversely, if the test passes with \emph{reasonable} probability,
$f$ must have a non-trivial agreement with another function $g$ that has the desired property.
The low degree test has been analyzed algebraically~\cite{FGLSS, AS, AroraSudan}
and also combinatorially~\cite{RazSafra}
whereas the linearity and dictatorship tests are often amenable to
Fourier analysis, \eg,~\cite{Kh-maxclique, Has96, Has97, KKMO}.

In this paper, we desire an explicit and \emph{good} trade-off between the alphabet size of the PCP and the
soundness parameter. At the Inner PCP level, this is achieved by a linearity test that combines
the elements of the linearity test and the low degree test. Specifically, given
a function  $f: \F_q^m \mapsto \F_q$, we wish to test that $f$ is linear. The standard
BLR test~\cite{BLR}
checks whether $f(x+y)=f(x)+f(y)$ for randomly chosen inputs $x$ and $y$. We
however wish to have a $2$-query test and hence we instead do a \emph{point versus subspace}
test. We are given the table of values of the function $f$ and in addition, for
every subspace $W \subseteq \F_q^m$ of dimension $6$, a linear function $\calT(W): W \mapsto \F_q$
that is
supposed to be the restriction of $f$ on $W$ (denoted $f|_W$). The test selects a random $6$-dimensional
subspace $W$ and a random point $w \in W$ and accepts if and only if
$f(w) = \calT(W) (w)$. Note that the query $f(w)$ is over an alphabet of size $q$ and the
query $\calT(W)$ is over an alphabet of size $q^6$.



We achieve the following soundness guarantee: if $f$ passes the \emph{point versus subspace}
test with probability ${3}/{q}$, then for some $j \in \F_q, j \not= 0$,
the function $j\cdot f$ has a Fourier coefficient with magnitude at least ${1}/{q^2}$
(see \expref{Lemma}{lemma:inner}). Interestingly, the test is analyzed by looking at the
probability that $f$ (and its restriction $f|_W$) passes the \emph{Gowers Test}
$f(x) - f(y) - f(z) + f(-x+y+z)=0$ for randomly chosen $x,y,z$ (see \expref{Section}{sec:gowers-test} for the reason the test is named so).
The Gowers Test is considered
for the purposes of the analysis only and is not a part of the actual test. The probability
of passing the Gowers Test is related to the existence of a \emph{large} Fourier coefficient
(see \expref{Lemma}{lemma:Gowers}) for function $j\cdot f$ for some $j \not= 0$.

The analysis proceeds as follows: assume that $f$ passes the \emph{point versus subspace}
test with probability ${1}/{q}+\delta$ where $\delta \geq {2}/{q}$. For the sake of
simplicity, assume that for every subspace $W$,
the test passes with probability ${1}/{q}+\delta$ after selecting $W$. Thus
$f|_W$ has agreement ${1}/{q}+\delta$ with a linear function $\calT(W)$.
We show that
this implies $j \cdot f|_W$ has a \emph{large} Fourier coefficient for some $j \not= 0$
and hence
$f|_W$ passes the Gowers Test with probability
${1}/{q}+\delta^4$. We observe that the probability that $f$ passes the Gowers Test is the average (over $W$)  of the probability that $f|_W$ passes the Gowers Test, up to an additive difference of
$e = {3}/{q^4}$. This implies that $f$ passes the Gowers Test with probability
${1}/{q} + \delta^4 - e \geq {1}/{q} + {2}/{q^4}$. {}From this we conclude that
for some $j \not= 0$, $j \cdot f$ has a \emph{large} Fourier coefficient, with magnitude
at least ${1}/{q^2}$.

Thus the Gowers Test serves as a vehicle to pass
from the local linearity of $f$ to its global linearity. This is also reminiscent of the
\emph{bootstrapping} method that allows the analysis of the low degree test in two (or three)
dimensions to carry over to a higher number of dimensions.

\paragraph{Remark} It is not clear that we necessarily need a \emph{point
versus $\ell$-dimensional subspace} test with $\ell=6$
to achieve a soundness of $O({1}/{q})$. Here is the
limitation of our current analysis.  In the Gowers Test
on $f|_W$, $dim(W)=\ell$, the
inputs $x, y, z$ are linearly dependent with probability
$\Theta({1}/{q^{\ell-2}})$. On the other hand, when the Gowers Test is
applied to the function $f$, the inputs $x,y,z$ are linearly dependent with probability
$\Theta({1}/{q^{m-2}} )$ which is negligible.  Hence we are able to claim that
the probability that $f$ passes the Gowers Test is the average (over $W$)  of the probability that $f|_W$ passes the Gowers Test, but only up to an additive difference of $e = \Theta({1}/{q^{\ell-2}})$. As the above calculation shows, we desire that $\delta
=O({1}/{q})$ and that $\delta^4$ dominates
$e = \Theta({1}/{q^{\ell-2}})$, which forces us to have $\ell \geq 6$.



\paragraph{Outer PCP, Composed PCP and Sub-Code Covering Property}
The linearity testing primitive at the Inner PCP level dictates that we use an Outer PCP
based on a \cclass{NP}-hard problem with linear constraints. A natural choice is the 3LIN problem
over $\F_q$: we are given an instance $(X, \Phi)$ where $X$ is a set of variables taking
values in $\F_q$ and $\Phi$ is a set of linear equations, each equation depending on three
variables from $X$. The goal is to find an assignment $\sigma: X \mapsto \Phi$ that
satisfies a good fraction of the equations. A well-known result of H\aa stad~\cite{Has97}
shows that for any constant $\eta > 0$,
it is \cclass{NP}-hard to distinguish whether an instance $(X,\Phi)$ has an assignment
that satisfies $1-\eta$ fraction of the equations (YES Case) or any assignment satisfies at most
${1}/{q}+\eta$ fraction of the equations (NO Case).

Starting with the hard instance of 3LIN as above, one builds a 2P1R Game as follows:
the first prover is sent a set of $k$ equations at random from $\Phi$. Let $V$ denote
the set of $3k$ variables sent to the first prover. The second prover is sent a set
$U \subseteq V$ that includes independently for $1 \leq i \leq k$, all three variables
in the $i^{th}$ equation with probability $1-\beta$ and exactly one of the three variables
in the $i^{th}$ equation with probability ${\beta}/{3}$ each
($\beta$ will be tiny as explained later).
The provers answer with assignments to $V$ and $U$ respectively and the verifier accepts
if and only if the assignments are consistent on $U$ and moreover all equations on $V$
are satisfied. In the YES Case, the provers have a strategy to make the verifier
accept with probability at least $1-k\eta$ whereas in the NO Case, any prover
strategy makes the verifier accept with probability at most $2^{-\Omega(\beta k)}$ (see
\expref{Section}{sec:outer} for a proof).

The 2P1R Game described is precisely the so-called Outer PCP\@. The composition of the
Outer and Inner PCP amounts to constructing a verifier that behaves as follows: the
PCP verifier expects, for each question $V$ ($U$ resp.) to the first (second resp.)
prover, a Hadamard encoding of assignment to $V$ ($U$ resp.). The Hadamard code is
same as the table of values of a linear function $f_V : \F_q^V \mapsto \F_q$ ($g_U: \F_q^U \mapsto \F_q$ resp.)
defined by an assignment to $V$ ($U$ resp.). Moreover, for every
$V$, a table of linear functions on all $6$-dimensional subspaces of $\F_q^V$ is
expected; these linear functions are supposed to be the restrictions of the global
linear function on $\F_q^V$. The verifier now picks a random question $V$ to the
first prover and performs the \emph{point versus subspace} test on the table
$f_V$ and the corresponding subspaces table.


Note that it appears as if the Hadamard codes on $U$ do not play any role (which would
not make sense). We observe that the Hadamard code on $U$ is actually contained in
the Hadamard code on $V$ (as a sub-code). This is because $\F_q^U \subseteq \F_q^V$ where
we append a vector in $\F_q^U$ with zeroes at positions in $V \setminus U$
and think of it as a vector in  $\F_q^V$. Thus there is no need to have a separate
Hadamard code on $U$ (though it helps in the analysis to think of these as virtual
tables). Moreover if $U \subseteq V \cap V'$ for distinct questions $V$ and $V'$ to the
first prover, we can identify the positions in Hadamard codes of $V$ and $V'$ that
correspond to the same position in the (virtual) Hadamard code of $U$.\footnote{The
reader might be concerned here that the 3LIN instance has imperfect
completeness and whether such identification can be done. This is a subtle and valid point:
a short answer is that
the identification can be done, but as mentioned in the next paragraph, the standard trick of \emph{folding over equations}
is no longer allowed and the \emph{folding} needs to be implemented explicitly.}

Now we look carefully at the soundness analysis of the composed PCP\@. Assume on the
contrary that the verifier
accepts with probability ${4}/{q}$ and for simplicity that for every $V$,
the verifier accepts
 the \emph{point versus subspace} test on the supposed Hadamard code
$f: \F_q^V \mapsto \F_q$ with probability ${4}/{q}$. The analysis of the Inner PCP
guarantees that for some $j \not=0$, $j\cdot f$ has a large Fourier coefficient.
Thus the function $f$ may be list decoded by making a list of all large Fourier
coefficients. Since the sum of squared magnitudes of all Fourier coefficients is $1$, the list size
is bounded. Our test also incorporates \emph{side-conditions} (see \expref{Section}{sec:point-subspace}) and ensures that
 the Fourier coefficients
obtained as list decoding satisfy all the equations on $V$ (this is done
in a more explicit manner than the standard \emph{folding over equations}
trick which is
inapplicable in our setting).

Finally we want to infer consistency between $f$ (\ie, supposed Hadamard code
on $V$) and the (virtual) supposed Hadamard codes $g_U: \F_q^U \mapsto \F_q$
on $U \subseteq V$. But as observed $g_U = f|_{\F_q^U}$.
We would like to conclude that since $j \cdot f$ has a large Fourier coefficient,
so do \emph{many}  of the $j \cdot g|_U$
functions. This turns out to be  possible if the sub-code spaces
$\F_q^U$ over all choices of $U$ (weighted according to the distribution on $U$ for
a fixed $V$) cover the global space $\F_q^V$ almost uniformly. We term this as
the \emph{sub-code covering property}.
When $\beta$ is sufficiently small and $k$ is
sufficiently large, we note that 
\[
|U| \approx \left(1-\frac{2}{3}\beta\right) |V|\,;
\]
thus the size of
a typical sub-code space is $q^{-O(\beta k)} = 2^{-O(\beta \log q \cdot k)}$
relative to the size of the code, there are roughly 
\[
\binom{k}{\beta  k} \cdot 3^{\beta k}
= 2^{\Omega(\beta \log(1/\beta)\cdot k)}
\]
choices of $U$, and it is
not unreasonable to expect that
the covering property holds provided $\log(1/\beta) \gg \log q$.
We formally prove this as \expref{Lemma}{lemma:covering}.

Once we are able to infer the consistency of tables $f = f_V$ and $g_U$,
as usual, the list decoding and then picking a random Fourier coefficient in the
list yields a provers' strategy in the 2P1R Game. This yields a contradiction
provided the soundness $2^{-\Omega(\beta k)}$ of the 2P1R Game
is low enough, which follows if
$\beta k$ is large enough.

\paragraph{The Codeword Test and the Consistency Test} In the PCP
literature, the low degree test, linearity test and the dictatorship test at the
Inner PCP level
are
often referred to as the \emph{codeword test} and then the Inner/Outer PCP composition
amounts to extending the test to the \emph{consistency test}
between two (or more as is necessary in some PCPs)
supposed codewords, \eg, $f_V$ and $g_U$ as above. Often this composition
presents serious technical challenges and one needs to carefully analyze
each PCP construction by itself. Our paper shows how to translate
the codeword test for Hadamard code (\ie, the linearity test), essentially in a
black-box manner, to a consistency test. In fact the tables $g_U$ are
virtual and there is no separate
consistency test. As described above, the codeword test for functions $f_V$
automatically serves as a consistency test between $f_V$ and $g_U$, since the
entries in tables
$f_V$ and  $f_{V'}$ that correspond to
the virtual table $g_U$ with $U \subseteq V \cap V'$
are identified together. We state this observation as an informal theorem.\footnote{We remark that a similar informal theorem also holds for the dictatorship test
modulo the Unique Games Conjecture, with the notion of Fourier coefficients replaced
by influences of co-ordinates.}

\section{Preliminaries}
In this section, we briefly describe preliminary background and the tools used in this paper.

\subsection{2-Prover 1-Round Games}

\begin{definition} \label{def:2p1r} A 2P1R Game ${\cal G}(\calV, \calU, \mu, {\cal R}, {\cal S},
\{ \pi_{VU}\} )$ consists of sets of questions $\calV, \calU$ and
sets of answers ${\cal R}, {\cal S}$ for the two provers
respectively, a distribution $\mu$ on the set of question pairs
$\calV \times \calU$ and for every question pair $(V,U)$ in the
support of $\mu$, a predicate $\pi_{VU}: {\cal R} \times {\cal S}
\mapsto \{0,1\}$ that defines the pairs of accepting answers.
A strategy of the %
provers is a map $\phi: \calV
\mapsto {\cal R}, \phi: \calU \mapsto {\cal S}$. The value of the
strategy $\phi$ is:
\[
  \val(\phi, {\cal G}) := \Pr_{(V,U) \sim \mu} \left[ \pi_{VU} (
  \phi(V), \phi(U) ) =1
 \right]\,.
\]
The value of the game ${\rm val}({\cal G})$ is the maximum value of
any prover strategy. A Projection Game is one where for every answer
of the first prover, there is exactly one accepting answer of the
second prover. For a Projection Game, the predicate $\pi_{VU}$ can
be thought of as a map $\pi_{VU}: {\cal R} \mapsto {\cal S}$ and
the accepting answers are of the form $(r, \pi_{VU}(r))$ for $r \in {\cal R}$.
For a Projection Game, $|{\cal S}| \leq |{\cal R}|$.
\end{definition}

A 2P1R Game is
best viewed as a game between the two provers and a verifier. The verifier
picks a random question pair $(V,U)$
from the distribution $\mu$, asks one question each to the two provers respectively,
and accepts if and only if the provers' answers satisfy the predicate $\pi_{VU}$.
The probability of acceptance
%
by
the verifier is same as the value of 
%
the
provers' strategy.

\begin{definition} Given a 2P1R Game ${\cal G}(\calV, \calU, \mu, {\cal R}, {\cal S},
\{ \pi_{VU}\} )$, the $k$-wise repeated game is
 \[
{\cal G}^{\otimes k}(\calV^k, \calU^k, \mu^k, {\cal R}^k, {\cal S}^k,
\{ \pi^k_{V'U'}\} )\,,
\] where for $V' = (V_1, \ldots, V_k)$ and $U' =
(U_1,\ldots, U_k)$, $\pi^k_{V'U'} := \bigwedge_{i=1}^k \pi_{V_iU_i}$.
\end{definition}

We state below Raz's Parallel Repetition Theorem along with the recent
improvements (and simplifications) by Holenstein and Rao.
\begin{theorem}[\cite{Raz, Holenstein, ARao}]\label{thm:pr} 
There exists an absolute constant $c > 0$ such that for a 2P1R
Game ${\cal G}$ with ${\rm val}({\cal G})= 1-\eps$,
 $$ {\rm val}({\cal G}^k) \leq (1-\eps^{3})^{c k}\,. $$
For a Projection Game, the bound of $(1-\eps^{2})^{c k}$ holds.
\end{theorem}

\subsection{Hardness of 3LIN}

Our reduction is from the 3LIN problem over a finite field. For the proof of \expref{Theorem}{thm:main}, we use H\aa stad's well-known hardness result for 3LIN~\cite{Has97}.
For the proof of Theorem~1.4, we need a hardness result for
3LIN with very good completeness and we use a result of Khot and Ponnuswami~\cite{KhotPonnuswami}.

\begin{definition} For a prime power $q$,
an instance $(X, \Phi)$ of 3LIN consists of a set of variables $X$
over $\F_q$ and a set of linear constraints $\Phi$ such that each
constraint depends on exactly three variables. Let $\OPT(X,
\Phi)$ denote the maximum fraction of constraints satisfied by any
assignment.
\end{definition}

\begin{theorem}[\cite{Has97}]\label{thm:Has97} 
For every constant $\eta > 0$ and a prime $q$, it is
\cclass{NP}-hard to distinguish whether a 3LIN instance $(X, \Phi)$ over
$\F_q$ has $\OPT(X,\Phi) \geq 1-\eta$ or $\OPT(X,\Phi)
\leq {1}/{q}+\eta$.
\end{theorem}

\begin{theorem}[\cite{KhotPonnuswami}]\label{thm:kp} 
There is a reduction from a SAT instance $\Gamma$
with size $n$ to a 3LIN instance $(X,\Phi)$ with size $N$ over a field $\F_q$ of characteristic $2$
such that: \begin{itemize}
\item $N$ and the running time of the reduction are bounded by $2^{O(\log^2 n)}$.
\item If $\Gamma$ is satisfiable, then $\displaystyle\OPT(X, \Phi) \geq
1- 2^{-\Omega(\sqrt{\log N})}$.
\item If $\Gamma$ is unsatisfiable, then $\displaystyle\OPT(X, \Phi) \leq
1- \Omega\left( \frac{1}{\log^3 N}  \right)$.
\end{itemize}
\end{theorem}

\begin{remark} \expref{Theorem}{thm:kp} is stated in~\cite{KhotPonnuswami} over the field $\F_2$. It also holds
for any extension field $\F_q$ since in the YES Case, the good assignment is over
$\F_2$ and in the NO Case, any assignment over $\F_q$ can be mapped
to an equally good assignment over $\F_2$, \eg, by taking the last bit of
every $\F_q$-element.
\end{remark}


\begin{remark} We present the proof of \expref{Theorem}{thm:main} by constructing a PCP over the prime field $\F_q$ and the Fourier analysis is over
this field. For the proof of Theorem~1.4, we use the same PCP but over
a field of characteristic $2$. We will omit the straightforward extension of our (Fourier) analysis to this case.
\end{remark}



\subsection{Hadamard code and Fourier analysis}

Let $q$ be a prime,   $\omega := e^{2\pi i/q}$ be the complex
$q^{th}$ root of unity and $\Omega := \{1, \omega, \ldots,
\omega^{q-1} \}$.

\begin{definition} 
%
%
The
%
Hadamard Code of $\alpha \in \F_q^n$ is defined
as the table of values of the linear function
 $\chi_\alpha: \F_q^n \mapsto \Omega$ where
  $$\forall x \in \F_q^n, \qquad  \chi_\alpha (x) = \omega^{\alpha \cdot  x}\,. $$
\end{definition}

The vector space of all functions $f : \F_q^n \mapsto {\mathbb C}$
has an orthonormal basis $\{ \chi_\alpha \mid \alpha \in \F_q^n\}$
where the inner product between two functions $f,g : \F_q^n \mapsto {\mathbb C}$ is defined as
(the expectation is over $x$ uniformly distributed in $\F_q^n$)
 $$ \left< f,g  \right> := \Exp_{x } \left[ f(x) \overline{g(x)}
 \right]\,. $$
Hence, every $f: \F_q^n \mapsto {\mathbb C}$ can be expressed
uniquely as
 $$ f = \sum_{\alpha \in \F_q^n}  \widehat{f}(\alpha) \ \chi_\alpha\,.
 $$
The coefficients $\widehat{f}(\alpha) \in {\mathbb C}$ are called
Fourier coefficients. These are defined by:
 $$ \widehat{f} (\alpha) = \left<f, \chi_\alpha \right> = \Exp_{x} \left[ f(x) \overline{\chi_\alpha(x)}
 \right]\,. $$
By Parseval's identity, $\sum_\alpha | \widehat{f}(\alpha)|^2
 = \| f\|_2^2 = \Exp_x [ |f(x)|^2]. $ In particular, for a function taking values in $\Omega$,
the sum of squared absolute values of all its Fourier coefficients equals $1$.


%
%
%
%



\subsection{Hellinger and statistical distance}

The squared Hellinger distance between distributions $D_1$ and $D_2$ over a discrete probability space
${\cal A}$ is
\[
H^2(D_1,D_2) := \frac{1}{2} \sum_{a\in {\cal A}}\left(\sqrt{D_1(a)}
  - \sqrt{D_2(a)}\right)^2 = 1 - \sum_{a\in {\cal A}}\sqrt{D_1(a)
  D_2(a)}\,.
\]
It is clear that $1-H^2(\cdot , \cdot)$ is multiplicative for product distributions $D_1^k, D_2^k$ on space
${\cal A}^k$, \ie,
  $$ 1-H^2(D_1^k, D_2^k) =  (1-H^2(D_1,D_2))^k\,.   $$
The statistical distance between $D_1$ and $D_2$ is:
$$\Delta(D_1,D_2) :=  \frac{1}{2} \sum_{a\in {\cal A}}\left|D_1(a) - D_2(a)\right|\,.$$
We have the standard inequality:
\begin{lemma}[\cite{Pol}]\label{p:Hellinger} 
$\displaystyle H^2(D_1,D_2) \leq \Delta(D_1,D_2) \leq \sqrt{2} \cdot H(D_1,D_2)\,.$
\end{lemma}









\section{The Outer PCP}\label{sec:outer}
In this section, we describe our Outer PCP\@. For the ease of exposition, we do it in three stages,
starting with the standard \emph{Variable Versus Equation} Game. We also prove the \emph{sub-code covering property}.
Let $(X, \Phi)$ be an instance of the 3LIN problem over $\F_q$ given by \expref{Theorem}{thm:Has97}
or \expref{Theorem}{thm:kp}.


\subsection{The Variable Versus Equation Game}

Consider the 2P1R Game where the verifier picks a random equation $E \in \Phi$ and then picks one of the three
variables $x \in E$. The first prover is sent the question $E$ (\ie, the three variables appearing in $E$)
and the second prover is sent the question $x$. The provers answer with the $\F_q$-values of all the variables
they receive.
The verifier accepts if and only if the two provers agree on the variable $x$ and moreover
the values given by the first prover
      satisfy the equation $E$. It is well-known and easy to check that
if $\OPT(X,\Phi) = 1-\eps$, then the value of the game is $1- {\eps}/{3}$.

\subsection{The Basic 2P1R Game}

We now slightly modify the Variable Versus Equation Game and call it the \emph{Basic 2P1R Game}
(think of $\beta$ as small):


\begin{itemize}
\item The verifier picks an equation $E \in \Phi$ at random. Let $x, y, z \in X$ be the variables in $E$.
\item The first prover is sent the equation $E$.
\item The second prover is sent the equation $E$ with probability $1-\beta$ and one of the three
variables $x,y, z$ with probability ${\beta}/{3}$ each.
\item The provers answer with the values of all the variables they receive.
\item The verifier accepts if and only if the two provers agree on the
      values of the variables and moreover the values given by the first prover
      satisfy the equation.
\end{itemize}

Assume that $\OPT(X,\Phi) = 1-\eps$.  The value of the game is at least $1-\eps$ since the provers
may stick to a $(1-\eps)$-satisfying assignment to $(X,\Phi)$ and in that case, the verifier may reject only
when the equation $E$ is not satisfied.

On the other hand, the value of the game is at most
$1-\Omega(\eps \beta)$ since with probability $\beta$, the second prover receives exactly one
variable and the variable versus equation game is played.

\subsection{The Final 2P1R Game (Outer PCP)}

The Outer PCP is now obtained as the $k$-wise repetition applied to the Basic 2P1R Game. Specifically:

\begin{itemize}
\item The verifier picks equations $E_1,\ldots,E_k \in \Phi$ at random. Let $V$ denote the set of $3k$ variables appearing in these equations.
\item The question $V$ is sent to the first prover.
\item Let $U \subseteq V$ be chosen by including independently for each $1 \leq i \leq k$, all three variables in equation $E_i$ with probability $1-\beta$ and one of the three variables in equation
    $E_i$ with probability ${\beta}/{3}$ each. The question $U$ is sent to the second prover.
\item The provers answer with the values of all the variables they receive.
\item The verifier accepts if and only if the two provers agree on the
      values of the variables in $U$ and moreover the values given by the first prover
      satisfy all the $k$ equations $E_1, \ldots, E_k$.
\end{itemize}

Assume again that $\OPT(X,\Phi) = 1-\eps$.  The value of the game is at least $1-\eps k $
since the provers
may stick to a $(1-\eps)$-satisfying assignment to $(X,\Phi)$ and in that case, the verifier may reject only
when at least one equation $E_i$ is not satisfied.

On the other hand, we observe that the value of the game is at most $(1-\Omega(\eps^2))^{\Omega(\beta k)}$.
As noted before, the value of the Basic Game is $1-\Omega(\eps \beta)$ and if we apply the parallel repetition
theorem (\expref{Theorem}{thm:pr}) directly, we end up with an upper bound of
$(1-\Omega(\eps^2 \beta^2))^{\Omega(k)}$. This bound is not good enough for us and
we get the better bound as follows.
Call a coordinate useful if the second prover receives a single variable, \ie,
the standard Variable Versus Equation Game is played. The expected number of useful
coordinates is $\beta k$. Hence the probability that less than $\beta k/2$
coordinates are useful is at most $2^{-\Omega(\beta k)}$ and may be ignored in comparison to the desired bound.
The repeated game restricted to the question pairs where at least $\beta k/2$
coordinates are useful can be thought of as a convex combination of sub-games,
each sub-game being the standard Variable Versus Equation Game repeated at least
$\beta k/2$ times. Each such sub-game has value at most $(1-\Omega(\eps^2))^{\Omega(\beta k)}$
and hence
so does the overall repeated game.

We note that the $k$ equations chosen by the verifier in the first step are
disjoint (\ie, do not share variables)
with high probability and we make this assumption in all that follows.


\subsection{Sub-code covering property}

Fix a question $(E_1, \ldots, E_k)$ to the first prover in the Outer PCP and
let $V$ denote the set of variables $\{x_1,\ldots,x_{3k}\}$ in these equations.
The question to the second prover is a subset $U \subseteq V$ where independently
for each $1 \leq i \leq k$,
all three variables
$x_{3(i-1)+1}, x_{3(i-1)+2}, x_{3i}$ are included in $U$ with probability $1-\beta$ and
exactly one of the three variables is included with probability ${\beta}/{3}$ each.
Consider two
distributions on $\F_q^V = \F_q^{3k}$:
\begin{itemize}
\item Distribution ${\cal D}$ is uniform on $\F_q^{V}$.
\item Distribution ${\cal D'}$
%
%
%
is obtained by picking
a random question $U \subseteq V$ to the second prover (as described above),
picking a string
in $\F_q^U$ uniformly at random, and then ``pulling up'' this string to
$\F_q^V$ by inserting $0$ at positions in $V \setminus U$.
\end{itemize}

\begin{lemma}\label{lemma:covering}  The statistical distance
$\Delta( {\cal D}, {\cal D'} )$ is upper bounded by $O(\sqrt{k}\cdot q\beta)$.
\end{lemma}
\begin{proof} We will bound the squared Hellinger distance on one coordinate, then
use the multiplicativity of the squared Hellinger distance for product distribution,
and then upper bound the statistical distance in terms of the Hellinger distance.

Let ${\cal Q}$ denote the uniform distribution on $\F_q$ and $({\cal Q}, {\cal Q}, {\cal Q})$
denote its three independent copies.
On any single coordinate $1 \leq j \leq k$,
note that ${\cal D}_j$ is same as the distribution $({\cal Q}, {\cal Q}, {\cal Q})$ whereas
${\cal D}'_j$ can be written as:
 $$ {\cal D}'_j = (1-\beta)\cdot ({\cal Q}, {\cal Q}, {\cal Q})+ \frac{\beta}{3} \cdot
 ({\cal Q}, 0, 0) + \frac{\beta}{3}\cdot (0, {\cal Q}, 0)+ \frac{\beta}{3} \cdot
 (0,0, {\cal Q})\,. $$
The total probability mass attached by ${\cal D}_j$ to triples in $\F_q^3$ where the
number of non-zero entries is three, two, one, and zero respectively is:
$$ p_3=\frac{(q-1)^3}{q^3}, \ p_2= \frac{3(q-1)^2}{q^3}, \ p_1=\frac{3(q-1)}{q^3}, \ p_0=\frac{1}{q^3}\,. $$
Similarly, the total probability mass attached by ${\cal D}'_j$ to triples in $\F_q^3$ where the
number of non-zero entries is three, two, one, and zero respectively is:
$$ p_3'=(1-\beta)\frac{(q-1)^3}{q^3}, \ p_2'=(1-\beta)\frac{3(q-1)^2}{q^3}, \ p_1'=(1-\beta)\frac{3(q-1)}{q^3}+
\beta \frac{q-1}{q}, \ p_0'=(1-\beta)\frac{1}{q^3}+ \beta \frac{1}{q}\,. $$
Hence
\begin{align*}
  1 - H^2({\cal D}_j, {\cal D}'_j )\
  &=\  \sqrt{p_3 p_3'} + \sqrt{p_2 p_2'} + \sqrt{p_1 p_1'} + \sqrt{p_0 p_0'} \\
  &=\  \frac{(q-1)^3}{q^3} \sqrt{1-\beta} + \frac{3(q-1)^2}{q^3}
  \sqrt{1-\beta}\\
  &\phantom{=\ \frac{(q-1)^3}{q^3} \sqrt{1-\beta}\ } +
  \frac{3(q-1)}{q^3} \sqrt{1+ \left(\frac{q^2}{3}-1\right)\beta} + \frac{1}{q^3} \sqrt{1+ (q^2-1)\beta} \\
  &\geq\ 1- O(\beta^2 q^2)\,,
\end{align*}
where all the square roots are estimated using
\[
\sqrt{1+x} \geq 1+\frac{x}{2}-x^2
\]
 for $x \in
[-{1}/{2},{1}/{2}]$
and $q^2 \beta \ll {1}/{2}$ (as will be the case). Thus the expression is
lower bounded by $a_0 + a_1 \beta + a_2 \beta^2$ where it is easily checked that
$a_0 =1 $ and $|a_2| = O(\beta^2 q^2)$.
The main point is that the coefficient $a_1$ vanishes. Indeed, $a_1$ equals
$$ \frac{1}{2} \left( - \frac{(q-1)^3}{q^3} -  \frac{3(q-1)^2}{q^3}
+  \frac{3(q-1)}{q^3} \cdot \left(\frac{q^2}{3}-1\right) +  \frac{1}{q^3} \cdot
  (q^2-1)  \right)   = 0\,.  $$
%
%
%
%
By multiplicativity of the squared Hellinger distance, we have
 $$ 1- H^2({\cal D}, {\cal D}') \ = \ (1- H^2({\cal D}_j, {\cal D}'_j))^{k}
 \ \geq \
 (1-O(\beta^2 q^2))^{k} \ \geq \ 1-O(\beta^2 q^2k)\,.$$
Finally $\Delta({\cal D}, {\cal D}') \leq  \sqrt{2} \cdot
 H({\cal D}, {\cal D}') \leq O(\sqrt{k} \cdot q\beta)$.
\end{proof}



\subsection{The choice of parameters}\label{sec:param}

Let $\OPT(X,\Phi)=1-\eps$ where $(X,\Phi)$ is a NO instance of 3LIN given by \expref{Theorem}{thm:Has97} or
\expref{Theorem}{thm:kp}.
\begin{itemize}
\item In \expref{Theorem}{thm:Has97}, $\OPT(X,\Phi)$ is close to ${1}/{q}$ and hence
$\eps = \Omega(1)$.
\item In \expref{Theorem}{thm:kp},
$\OPT(X,\Phi) \leq 1-\Omega({1}/{\log^3 N})$ where $N$ is the
size of the 3LIN instance and hence $\eps = \eps(N) =  \Omega({1}/{\log^3 N})$.
\end{itemize}
The parameters will be chosen so that for a large enough constant $C$,
\begin{itemize}
\item The soundness of the Outer PCP, which is at most $(1-\Omega(\eps^2))^{-\Omega(\beta k)}$,
      is at most ${1}/{q^{C}}$.
\item $\displaystyle \Delta({\cal D}, {\cal D}') \leq \frac{1}{C q^2}$.
\end{itemize}
Using \expref{Lemma}{lemma:covering},
it suffices to choose 
\[
k = \frac{C_*^3  q^6 \log^2 q}{\eps^4} \quad \text{and}\quad
\beta = \frac{\eps^2}{C_*^2 \cdot q^6 \cdot \log q}
\]
for a large enough constant $C_*$.







%
%
%
%
%
%
%
%
%
%
%
%
%
%



\section{The Inner PCP}

In this section, we describe our Inner PCP\@. We first analyze a test that we call Gowers Test,
which is then used to analyze the actual Inner PCP presented in \expref{Section}{sec:point-subspace}. We begin with
some notation and a simple lemma.

Let $\omega := e^{2\pi i/q}$ be the complex $q^{th}$ root of unity and $\Omega := \{1, \omega,
\ldots, \omega^{q-1} \}$. For $f, g : \F_q^m \mapsto \Omega$, let $\AGR(f,g)$ denote the agreement
between the two functions, \ie, the fraction of points on which they agree.
Let $f_i$ denote the function $f_i(x) := f(x)^i$. Note that for $z \in \Omega$, the expression
$(1+z+z^2+\cdots+z^{q-1})/q$ equals $1$ if $z=1$ and $0$ otherwise.

\begin{lemma}\label{lemma:agr} If $f: \F_q^m \mapsto \Omega$ is a function such that for some
linear function $\chi_\alpha$, $\AGR(f, \chi_\alpha) \geq {1}/{q}+\delta$. Then
$\sum_{j=1}^{q-1} \left| \widehat{f}_{j}(j \alpha) \right| \geq q\delta$.
\end{lemma}
\begin{proof} The lemma follows by noting that:
\begin{align*}
\AGR(f, \chi_\alpha)\ &=\  \Exp_x \left[  \frac{1}{q} \sum_{j=0}^{q-1} (f(x) \overline{\chi_\alpha(x)})^j   \right] \\
 &=\   \frac{1}{q} + \frac{1}{q} \sum_{j=1}^{q-1} \Exp_x \left[ f_j(x) \overline{\chi_{j \alpha}(x)}   \right] \\
 &=\  \frac{1}{q} + \frac{1}{q}  \sum_{j=1}^{q-1} \widehat{f}_j (j \alpha)\,.\qedhere
\end{align*}
\end{proof}

\subsection{The Gowers test}\label{sec:gowers-test}

For a function $f: \F_q^m \mapsto {\mathbb C}$, the Gowers Uniformity Norm $U_2$ \cite{Gowers} is defined as
 $$ \| U_2(f) \|^4 := \Exp_{x,y,z} \left[ f(x) \overline{f(x+y)}\overline{f(x+z)} f(x+y+z)   \right]\,. $$
We will study the probability that a function $f: \F_q^m \mapsto \Omega$ passes the test
$$f(x) \overline{f(x+y)}\overline{f(x+z)} f(x+y+z) =1\,,$$
for a random choice of $x,y,z$. It is thus natural to name the test as the Gowers Test.
It will be more convenient for us to think of the test equivalently
as $$\mbox{Gowers Test}: \quad\quad f(x) \overline{f(y)}\overline{f(z)} f(-x+y+z) =1\,.$$

\begin{lemma} Let $f: \F_q^m \mapsto \Omega$ be a function. Then the acceptance probability of the Gowers Test is:
$$ \Pr_{x,y,z} \left[f(x) \overline{f(y)}\overline{f(z)} f(-x+y+z) =1 \right]
   = \frac{1}{q} + \frac{1}{q} \sum_{j=1}^{q-1} \sum_\alpha  \left| \widehat{f}_j (j\alpha)\right|^4\,.   $$
\end{lemma}
\begin{proof} The acceptance probability can be expressed as
\begin{align*}
   &  \frac{1}{q} + \frac{1}{q} \sum_{j=1}^{q-1} \Exp \left[ (f(x) \overline{f(y)}\overline{f(z)} f(-x+y+z))^j   \right] \\
  =\ & \frac{1}{q} + \frac{1}{q}  \sum_{j=1}^{q-1} \sum_{\alpha, \phi, \psi, \gamma}
 \widehat{f_j}(\alpha) \overline{\widehat{f_j}(\phi)} \overline{\widehat{f_j}(\psi)}\widehat{f_j}(\gamma)
  \Exp \left[ \chi_{\alpha - \gamma}(x) \chi_{-\phi + \gamma}(y) \chi_{-\psi + \gamma}(z)    \right] \\
  =\ &  \frac{1}{q} + \frac{1}{q}  \sum_{j=1}^{q-1} \sum_\alpha \left| \widehat{f}_j (\alpha)\right|^4\,,
\end{align*}
noting that the expectation vanishes unless $\alpha = \phi = \psi = \gamma$. We can replace
$\alpha$ by $j \alpha$ without changing the summation.
\end{proof}

\subsection{The Gowers test with side conditions}

At the Inner PCP level, we need to check not only that a function $f$ is linear, \ie, $f= \chi_\alpha$ for some
$\alpha$, but also that $\alpha$ itself satisfies a given set of linear constraints. We call these
as \emph{side conditions} and modify the Gowers Test so as to incorporate these side conditions (and henceforth Gowers Test refers to one with side conditions incorporated).

\medskip
\noindent {\bf Given:}
\begin{itemize}
\item A function $f: \F_q^m \mapsto \Omega$.
\item Side conditions $h_i \cdot x = b_i, i = 1,\ldots, k$ where $h_i \in \F_q^m, \ b_i \in \F_q$. Assume that $\{h_i\}_{i=1}^k$ are linearly independent and  $H$ is their
    linear span.
\end{itemize}
\noindent {\bf The Test:}
\begin{itemize}
\item Pick $x, y, z \in \F_q^m$ at random.
\item Pick $a = (a_1,\ldots, a_k) \in \F_q^k$ at random and let $h= \sum_{i=1}^k a_i h_i$ and
      $a\cdot b := \sum_{i=1}^k a_i b_i$.
\item Accept if and only if
  $$f(x) \overline{f(y)}\overline{f(z)} f(-x+y+z+h) = \omega^{a\cdot b}\,. $$
\end{itemize}


\begin{lemma}\label{lemma:Gowers} The following hold:
\begin{enumerate}
\item The Gowers Test always passes with probability at least ${1}/{q}$.
\item If the Gowers Test passes with probability at least ${1}/{q}+\delta$, then there
exists $1 \leq j \leq q-1$ and $\alpha \in \F_q^m$ that respects the side conditions (\ie,
$\forall i, \alpha \cdot h_i = b_i$) and $| \widehat{f}_j(j \alpha) | \geq \sqrt{\delta}$.
\item If $f$ has agreement ${1}/{q}+\delta$ with a linear function $\chi_\alpha$ that respects
the side conditions (\ie,
$\forall i, \alpha \cdot h_i = b_i$), then $f$ passes the Gowers Test with probability at least
${1}/{q}+ {\delta^4}$.
\end{enumerate}
\end{lemma}
\begin{proof} The acceptance probability of the Gowers Test with side conditions is:
\begin{align*}
   &  \frac{1}{q} + \frac{1}{q} \sum_{j=1}^{q-1} \Exp \left[ (f(x) \overline{f(y)}\overline{f(z)} f(-x+y+z+h))^j   \omega^{-j a\cdot b} \right] \\
  =\ & \frac{1}{q} + \frac{1}{q}  \sum_{j=1}^{q-1} \sum_{\alpha, \phi, \psi, \gamma}
 \widehat{f_j}(\alpha) \overline{\widehat{f_j}(\phi)} \overline{\widehat{f_j}(\psi)}\widehat{f_j}(\gamma)
  \Exp \left[ \chi_{\alpha - \gamma}(x) \chi_{-\phi + \gamma}(y) \chi_{-\psi + \gamma}(z)    \right]
  \Exp_a \left[ \chi_{\gamma}(h) \omega^{-j a\cdot b}    \right]  \\
  =\ &  \frac{1}{q} + \frac{1}{q}  \sum_{j=1}^{q-1} \sum_\alpha \left| \widehat{f}_j (\alpha)\right|^4
  \cdot \Exp_a \left[ \omega^{\sum_{i=1}^k a_i (\alpha \cdot h_i - j b_i)}  \right] \\
  =\ &  \frac{1}{q} + \frac{1}{q}  \sum_{j=1}^{q-1} \ \  \sum_{\alpha| \forall i, \alpha\cdot h_i = j b_i}
  \ \ \left| \widehat{f}_j (\alpha)\right|^4 \\
   =\ &  \frac{1}{q} + \frac{1}{q}  \sum_{j=1}^{q-1} \ \  \sum_{\alpha| \forall i, \alpha\cdot h_i =  b_i}
  \ \ \left| \widehat{f}_j (j\alpha)\right|^4\,. \\
\end{align*}
The first conclusion is immediate. The second uses the fact that the sum of squared
absolute values of Fourier coefficients of $f_j$ is $1$.
 The third follows in conjunction with \expref{Lemma}{lemma:agr} and convexity of the
 function $x^4$.
\end{proof}

\subsection{The Point-Subspace test with side conditions (Inner PCP)}\label{sec:point-subspace}

\noindent {\bf Given:}
\begin{itemize}
\item A function $f: \F_q^n \mapsto \Omega$.
\item Side conditions $h_i \cdot x = b_i, i =1,\ldots, k$. Assume that $\{h_i\}_{i=1}^k$ are linearly independent and $H$ is their
    linear span.
\item A table $\left\{ {\cal T}(W) \mid  W \in {\cal C} \right\}$ where ${\cal C}$ denotes the class of $k+6$ dimensional subspaces of the form $W = D \oplus H$ for a $6$-dimensional subspace $D$ of $\F_q^n$ such that $D \cap H = \{ {\bf 0} \}$. The entry ${\cal T}(W)$ is a linear function $\chi: W \mapsto \Omega$ that respects the side conditions, \ie,
\[
\chi(x+y) = \chi(x)\cdot \chi(y)\qquad\text{and}\qquad \chi(x + h_i) =
\chi(x)\cdot  \omega^{b_i} \quad \forall x,y \in W, i=1,\ldots,k\,.
\]
\end{itemize}
\noindent {\bf The Test:}
\begin{enumerate}
\item Pick a random $W \in {\cal C}$ and let $ {\cal T}(W)$ be the linear function on $W$.
\item Pick a random $w \in W$.
\item Accept if and only if $f(w) = {\cal T}(W)(w)$.
\end{enumerate}

\subsubsection{Completeness}
Suppose $f: \F_q^n \mapsto \Omega$ is linear that respects the side
conditions, \ie, \[
f(x+y) = f(x)\cdot f(y) \qquad\text{and}\qquad f(x + h_i) = f(x)\cdot
\omega^{b_i} \quad \forall x,y \in \F_q^n, i=1,\ldots,k\,.
\]
Then letting
${\cal T}(W)$ to be the restriction $f|_W$, the point-subspace test passes with probability $1$.

\subsubsection{Soundness}

\begin{lemma} \label{lemma:inner} If the point-subspace test passes with probability
${3}/{q}$ then there exists
$1 \leq j \leq q-1$ and $\alpha\in \F_q^n$ that respects the side conditions (\ie, $\forall i,
\alpha \cdot h_i = b_i$) and $| \widehat{f}_j (j \alpha)| \geq {1}/{q^2}$.
\end{lemma}
\begin{proof} Assume that the point-subspace test passes
with probability ${1}/{q}+\delta$ where $\delta \geq {2}/{q}$. This means that:
  $$ \Exp_W \left[ \AGR\left( f|_W,  {\cal T}(W) \right) \right]  = {1}/{q}+\delta\,. $$
By \expref{Lemma}{lemma:Gowers}(1,3) and using convexity of the function $x^4$,
we conclude that
\begin{equation}
\Exp_W \left[\Pr \left[ f|_W \ \mbox{passes Gowers Test} \right] \right] \geq {1}/{q}+\delta^4.
\end{equation}
Now let us consider the probability that $f$ passes the Gowers Test. The Gowers Test
picks three points $x, y, z$ independently (from the global space $\F_q^n$). Let
$\bot$ be the event that  ${\rm span}(x,y,z,H)$ has dimension $k+3$ and
let
$\Gamma$ be the set of all such triples $(x,y,z)$. Conditional on $\bot$ happening,
the Gowers Test picks a triple in $\Gamma$ uniformly at random. An alternate way of
picking a triple in $\Gamma$ uniformly at random is to pick a subspace $W \in {\cal C}$
and then pick $(x,y,z) \in W^3$
conditional on the event that ${\rm span}(x,y,z,H)$ has dimension $k+3$. Let $\bot(W)$
be the event that ${\rm span}(x,y,z,H)$ has dimension $k+3$
when $x,y,z \in W$ are picked at random. Thus:
\begin{align*}
\Pr \left[ \text{$f$ passes Gowers Test} \mid \bot \right]\ &
 =  \ \Exp_W \bigl[ \Pr\left[ f|_W \ \mbox{passes Gowers Test} \mid \bot(W) \right]   \bigr]     \\[1ex]
 & \geq\   \Exp_W \bigl[ \Pr \left[ f|_W \ \mbox{passes Gowers Test}\right] - \Pr[\neg \bot(W)] \bigr] \\
 & \geq\   \Exp_W \bigl[ \Pr \left[ f|_W \ \mbox{passes Gowers Test}\right] \bigr]  - \frac{3}{q^4} \\
 & \geq\   \frac{1}{q} +\delta^4 -\frac{3}{q^4} \ \geq \ \frac{1}{q}+\frac{2}{q^4}\,,
 \end{align*}
where we used $\Pr[ \neg \bot(W)] \leq {3}/{q^4}$ and $\delta \geq {2}/{q}$.
Also, noting that 
\[
\Pr[\bot] \geq 1-\frac{3}{q^{n-k-2}}\,,
\]
we get that
$$\Pr\left[f \ \mbox{passes Gowers Test}\right] \geq \frac{1}{q}+\frac{1}{q^4}\,.$$
It follows now from \expref{Lemma}{lemma:Gowers}(2)
that there exists $1 \leq j \leq q-1$ and $\alpha$ such that $\forall i,
\alpha \cdot h_i = b_i$ and $| \widehat{f}_j (j \alpha)| \geq {1}/{q^2}$.
\end{proof}


\section{The composed PCP}\label{sec:composed}
We now describe the composed PCP and prove \expref{Theorem}{thm:main}.
The Outer PCP is constructed from a 3LIN instance $(X, \Phi)$
as described in \expref{Section}{sec:outer}. The 3LIN instance is either $(1-\eta)$-satisfiable
or at most $(1-\eps)$-satisfiable as per \expref{Theorem}{thm:Has97} or \expref{Theorem}{thm:kp}.
The various parameters are chosen as in \expref{Section}{sec:param}.


The verifier in the composed PCP expects the proof
to contain, for every question $V$ to the first prover, two tables $\calL_V$
and ${\cal T}_V$.

\paragraph{The tables $\calL_V$}
The table $\calL_V$ gives
the Hadamard code of the
assignment to $V$. Concretely, let $\{x_1,\ldots,x_{3k}\}$ be the variables in $V$
and $\sigma : X \mapsto \F_q$ be the global assignment that is supposed to be an
almost satisfying assignment to $(X,\Phi)$. The verifier expects, for the question $V$,
the table of values of the linear function $\calL_V: \F_q^V = \F_q^{3k}
\mapsto \Omega$ where
 $$ \calL_V(y) = \omega^{\sum_{i=1}^{3k} y_i \sigma(x_i)}\,. $$
Note that for every question $U$ to the second prover, a table $\calL_U$ that gives the
Hadamard code
of the assignment to $U$ may be expected as well. However if $U \subseteq V$, then the table $\calL_U$
is contained in the table $\calL_V$. Specifically, for any $z \in \F_q^U$,
let $z^{\uparrow}$ denote the vector obtained by extending $z$
to a vector in $\F_q^V$ be inserting
zeroes at the positions in $V\setminus U$. Then
 $$ \calL_U (z) = \omega^{\sum_{i: x_i \in U} z_i \sigma(x_i)} = \omega^{\sum_{i=1}^{3k}
    z^{\uparrow}_i \sigma(x_i)}
  = \calL_V (z^{\uparrow})\,. $$
Thus there is no need to have separate $\calL_U$ tables; we do however
think of these as virtual tables. Whenever
there are questions $V, V'$ to the first prover such that $U \subseteq V \cap V'$,
the table $\calL_U$ is contained in both the tables $\calL_V$ and $\calL_{V'}$.
We identify the
locations in these two tables that correspond to the same location in the
(virtual) $\calL_U$ table.

\paragraph{The tables $\calT_V$}
Let $V$ be a question to the first prover that consists of equations (\ie,
side conditions)
$\{h_i\cdot x = b_i\}_{i=1}^k$, the $i^{th}$ equation depending only on the
variables $(x_{3(i-1)+1}, x_{3(i-1)+2}, x_{3i})$. Except with a negligible
probability (that we may ignore) over
the choice of $V$, the $k$ equations are disjoint (\ie, do not share variables).
Thus the vectors $h_i$ are
linearly independent and let $H$ be their linear span.
The table $\calT_V$ contains, for every $k+6$ dimensional subspace $W \subseteq
\F_q^V$ such that $H\subseteq W$, a linear function $\calT_V(W) : W \mapsto \Omega$
that respects the side conditions.

\paragraph{The PCP verifier}  The verifier picks a random question
$V$ to the first prover in the Outer PCP\@. Let $\{h_i\cdot x = b_i\}_{i=1}^k$ be the side conditions and $H$ be the linear span of vectors $\{h_i\}_{i=1}^k$. Let
$\calL_V: \F_q^V \mapsto \Omega$ and $\calT_V$ be the
associated tables. The verifier picks a random $k+6$ dimensional
subspace $W$, $H \subseteq W \subseteq \F_q^V$, and a random point $w \in W$.
The verifier accepts if and only if
 $$  \calL_V(w) = \calT_V(W) (w)\,. $$

\subsection{Completeness}

Let $\sigma: X \mapsto \F_q$ be a global assignment that satisfies $1-\eta$ fraction
of equations. The table $\calL_V$ is the Hadamard code (\ie, linear function)
of the assignment $\sigma$ restricted
to $V$. The table $\calT_V$ gives, for every subspace $W$, the linear function
$\calT_V(W) = \calL_V|_W$.
The test may fail only when there is some equation in $V$ that is not satisfied
by $\sigma$. This happens with probability at most $\eta k$.



\subsection{Soundness}

Assume on the contrary that the test accepts with probability ${4}/{q}$. By an averaging argument,
for at least ${1}/{q}$ fraction of the questions $V$, the test accepts
with probability at least ${3}/{q}$. Fix any such \emph{good} $V$.
By the analysis of the Inner PCP, \expref{Lemma}{lemma:inner},
it follows that there exists $1 \leq j\leq q-1$ such that for
$f := \calL_V$, there is a Fourier coefficient $|\widehat{f_j}(j\alpha)| \geq {1}/{q^2}$ and
$\alpha$ respects the side conditions.
Note that for the uniform distribution ${\cal D}$ on $\F_q^V$,
 $$ \widehat{f_j}(j\alpha) = \Exp_{x \in {\cal D}} \left[ f_j(x)  \overline{\chi_{j\alpha}(x)}   \right]\,. $$
By the sub-code covering property, \expref{Section}{sec:param}, it follows that for some error parameter
$e, |e| \leq {1}/{(C q^2)}$,
\begin{align*}
\widehat{f_j}(j\alpha)\ &= \ \Exp_{x \in {\cal D'}} \left[  f_j(x) \overline{\chi_{j\alpha}(x)}   \right] \  + e \\
 & = \  \Exp_U \left[   \Exp_{y \in \F_q^U} \left[ f_j(y^\uparrow)
 \overline{\chi_{j\alpha}(y^\uparrow)}  \right]  \ \right] \ + e \\
 & = \  \Exp_U \left[   \Exp_{y \in \F_q^U} \left[ f_j|_U(y)
 \overline{\chi_{j\alpha_\downarrow}(y)}  \right]  \ \right] \ + e \\
 & = \  \Exp_U \left[ \widehat{f_j|_U} (j \alpha_\downarrow) \right] \ + e\,,
\end{align*}
where
$f_j|_U$ denotes the restriction of $f_j$ to $\F_q^U \subseteq \F_q^V$
and $\alpha_\downarrow$
denotes the vector obtaining by dropping coordinates of $\alpha$ in $V \setminus U$.
It follows that
 $$  \Exp_U \left[ \ \left| \widehat{f_j|_U} (j \alpha_\downarrow) \right| \right]
\geq |\widehat{f_j}(j\alpha)| - |e| \geq \frac{1}{q^2}- \frac{1}{Cq^2} \geq \frac{1}{2q^2}\,. $$
Thus, with probability at least
${1}/{(4q^2)}$ over the choice of $U$ (for the fixed good $V$; call such $(V,U)$  a \emph{good} pair),
\[
\Bigl|\widehat{f_j|_U}({j \alpha_\downarrow})\Bigr| \geq \frac{1}{4q^2}\,.
\]
Note also that
$f|_U = g = \calL_U$ which is the supposed (virtual) Hadamard code for an assignment to $U$ and
$f_j|_U = g_j$.

%
Now we derive strategies for the provers in the Outer PCP as follows: the first prover,
on receiving a question $V$, considers the function $f = \calL_V$,
lists all $\alpha$ such that for some $1 \leq j \leq q-1$, $|\widehat{f_j}(j\alpha)| \geq {1}/{q^2}$ and $\alpha$ respects the side conditions, and
then outputs a random element from the list. The list size is bounded by $q^5$.
The second prover, on receiving a question $U$, considers the function
$g = \calL_U$,
lists all $\gamma$
such that for some $1 \leq j \leq q-1$, $|\widehat{g_j}(j\gamma)| \geq {1}/{(4q^2)}$, and
outputs a random element from the list. The list size is bounded by $16q^5$.
The analysis above shows that when $(V,U)$ is a good pair, there are $\alpha$
and $\alpha_\downarrow = \gamma$ in the lists for $V$ and $U$ respectively and $\alpha$
respects the side conditions. This and the bound on the list sizes
shows that with probability at least 
\[
\frac{1}{q}\cdot \frac{1}{4q^2} \cdot\frac{1}{q^5}
\cdot \frac{1}{16q^5}\,,
\]
the provers succeed. This is a contradiction since the
soundness of the Outer PCP is at most ${1}/{q^C}$ for a large constant $C$ as in
\expref{Section}{sec:param}.

%
%
%
%
%
\providecommand{\bibhead}[1]{}
%
\expandafter\ifx\csname pdfbookmark\endcsname\relax%
  \providecommand{\tocrefpdfbookmark}{}
\else\providecommand{\tocrefpdfbookmark}{%
   \hypertarget{tocreferences}{}%
   \pdfbookmark[1]{References}{tocreferences}}%
\fi

\tocrefpdfbookmark
\begin{thebibliography}{10}

\bibitem{ABSS}\bibhead{ABSS}
{\sc Sanjeev Arora, L{\'a}szl{\'o} Babai, Jacques Stern, and Elizabeth~Z.
  Sweedyk}: The hardness of approximate optima in lattices, codes, and systems
  of linear equations.
\newblock {\em J. Comput. System Sci.}, 54(2):317--331, 1997.
\newblock Preliminary version in
  \href{http://dx.doi.org/10.1109/SFCS.1993.366815}{FOCS'93}.
\newblock [\epfmtdoi{10.1006/jcss.1997.1472}]

\bibitem{ABHKS}\bibhead{ABHKS}
{\sc Sanjeev Arora, Eli Berger, Elad Hazan, Guy Kindler, and Muli Safra}: On
  non-approximability for quadratic programs.
\newblock In {\em Proc. 46th FOCS}, pp. 206--215. IEEE Comp. Soc. Press, 2005.
\newblock See also at \href{http://eccc.hpi-web.de/report/2005/058/}{ECCC}.
\newblock [\epfmtdoi{10.1109/SFCS.2005.57}]

\bibitem{ALMSS}\bibhead{ALMSS}
{\sc Sanjeev Arora, Carsten Lund, Rajeev Motwani, Madhu Sudan, and Mario
  Szegedy}: Proof verification and the hardness of approximation problems.
\newblock {\em J. ACM}, 45(3):501--555, 1998.
\newblock Preliminary version in
  \href{http://dx.doi.org/10.1109/SFCS.1992.267823}{FOCS'92}. See also at
  \href{http://eccc.hpi-web.de/report/1998/008/}{ECCC}.
\newblock [\epfmtdoi{10.1145/278298.278306}]

\bibitem{AS}\bibhead{AS}
{\sc Sanjeev Arora and Shmuel Safra}: Probabilistic checking of proofs: A new
  characterization of $\mathsf{NP}$.
\newblock {\em J. ACM}, 45(1):70--122, 1998.
\newblock Preliminary version in
  \href{http://dx.doi.org/10.1109/SFCS.1992.267824}{FOCS'92}.
\newblock [\epfmtdoi{10.1145/273865.273901}]

\bibitem{AroraSudan}\bibhead{AroraSudan}
{\sc Sanjeev Arora and Madhu Sudan}: Improved low-degree testing and its
  applications.
\newblock {\em Combinatorica}, 23(3):365--426, 2003.
\newblock Preliminary version in
  \href{http://dx.doi.org/10.1145/258533.258642}{STOC'97}.
\newblock [\epfmtdoi{10.1007/s00493-003-0025-0}]

\bibitem{BGS}\bibhead{BGS}
{\sc Mihir Bellare, Oded Goldreich, and Madhu Sudan}: Free bits, {PCPs}, and
  nonapproximability---towards tight results.
\newblock {\em SIAM J. Comput.}, 27(3):804--915, 1998.
\newblock Preliminary version in
  \href{http://dx.doi.org/10.1109/SFCS.1995.492573}{FOCS'95}. See also at
  \href{http://eccc.hpi-web.de/report/1995/024/}{ECCC}.
\newblock [\epfmtdoi{10.1137/S0097539796302531}]

\bibitem{BGLR}\bibhead{BGLR}
{\sc Mihir Bellare, Shafi Goldwasser, Carsten Lund, and A.~Russell}: Efficient
  probabilistically checkable proofs and applications to approximations.
\newblock In {\em Proc. 25th STOC}, pp. 294--304. ACM Press, 1993.
\newblock See correction to abstract in
  \href{http://dx.doi.org/10.1145/195058.195467}{STOC'94}.
\newblock [\epfmtdoi{10.1145/167088.167174}]

\bibitem{BLR}\bibhead{BLR}
{\sc Manuel Blum, Michael Luby, and Ronitt Rubinfeld}: Self-testing/correcting
  with applications to numerical problems.
\newblock {\em J. Comput. System Sci.}, 47(3):549--595, 1993.
\newblock Preliminary version in
  \href{http://dx.doi.org/10.1145/100216.100225}{STOC'90}.
\newblock [\epfmtdoi{10.1016/0022-0000(93)90044-W}]

\bibitem{Chan}\bibhead{Chan}
{\sc Siu~On Chan}: Approximation resistance from pairwise independent
  subgroups.
\newblock In {\em Proc. 45th STOC}, pp. 447--456. ACM Press, 2013.
\newblock See also at \href{http://eccc.hpi-web.de/report/2012/110/}{ECCC}.
\newblock [\epfmtdoi{10.1145/2488608.2488665}]

\bibitem{CharikarWirth}\bibhead{CharikarWirth}
{\sc Moses Charikar and Anthony Wirth}: Maximizing quadratic programs:
  extending {G}rothendieck's inequality.
\newblock In {\em Proc. 45th FOCS}, pp. 54--60. IEEE Comp. Soc. Press, 2004.
\newblock [\epfmtdoi{10.1109/FOCS.2004.39}]

\bibitem{condon}\bibhead{condon}
{\sc Anne Condon}: The complexity of the max word problem and the power of
  one-way interactive proof systems.
\newblock {\em Comput. Complexity}, 3(3):292--305, 1993.
\newblock Preliminary version in
  \href{http://dx.doi.org/10.1007/BFb0020820}{STACS'91}.
\newblock [\epfmtdoi{10.1007/BF01271372}]

\bibitem{EH-kCSP}\bibhead{EH-kCSP}
{\sc Lars Engebretsen and Jonas Holmerin}: More efficient queries in {PCPs} for
  $\mathsf{NP}$ and improved approximation hardness of maximum {CSP}.
\newblock {\em Random Structures {\&} Algorithms}, 33(4):497--514, 2008.
\newblock Preliminary version in
  \href{http://dx.doi.org/10.1007/978-3-540-31856-9_16}{STACS'05}.
\newblock [\epfmtdoi{10.1002/rsa.20226}]

\bibitem{Feige}\bibhead{Feige}
{\sc Uriel Feige}: A threshold of $\ln n$ for approximating set cover.
\newblock {\em J. ACM}, 45(4):634--652, 1998.
\newblock Preliminary version in
  \href{http://dx.doi.org/10.1145/237814.237977}{STOC'96}.
\newblock [\epfmtdoi{10.1145/285055.285059}]

\bibitem{FGLSS}\bibhead{FGLSS}
{\sc Uriel Feige, Shafi Goldwasser, L{\'a}szl{\'o} Lov{\'a}sz, Shmuel Safra,
  and Mario Szegedy}: Interactive proofs and the hardness of approximating
  cliques.
\newblock {\em J. ACM}, 43(2):268--292, 1996.
\newblock Preliminary version in
  \href{http://dx.doi.org/10.1109/SFCS.1991.185341}{FOCS'91}.
\newblock [\epfmtdoi{10.1145/226643.226652}]

\bibitem{Gowers}\bibhead{Gowers}
{\sc Timothy Gowers}: A new proof of {S}zemer{\'e}di's theorem for progressions
  of length four.
\newblock {\em Geometric and Functional Analysis}, 8(3):529--551, 1998.
\newblock [\epfmtdoi{10.1007/s000390050065}]

\bibitem{Has96}\bibhead{Has96}
{\sc Johan H{\aa}stad}: Clique is hard to approximate within $n^{1-\epsilon}$.
\newblock {\em Acta Mathematica}, 182(1):105--142, 1999.
\newblock Preliminary version in
  \href{http://dx.doi.org/10.1109/SFCS.1996.548522}{FOCS'96}.
\newblock [\epfmtdoi{10.1007/BF02392825}]

\bibitem{Has97}\bibhead{Has97}
{\sc Johan H{\aa}stad}: Some optimal inapproximability results.
\newblock {\em J. ACM}, 48(4):798--859, 2001.
\newblock Preliminary version in
  \href{http://dx.doi.org/10.1145/258533.258536}{STOC'97}.
\newblock [\epfmtdoi{10.1145/502090.502098}]

\bibitem{Holenstein}\bibhead{Holenstein}
{\sc Thomas Holenstein}: Parallel repetition: Simplification and the
  no-signaling case.
\newblock {\em Theory of Computing}, 5(8):141--172, 2009.
\newblock Preliminary version in
  \href{http://dx.doi.org/10.1145/1250790.1250852}{STOC'07}.
\newblock [\epfmtdoi{10.4086/toc.2009.v005a008}]

\bibitem{Kh-maxclique}\bibhead{Kh-maxclique}
{\sc Subhash Khot}: Improved inapproximability results for {MaxClique},
  chromatic number and approximate graph coloring.
\newblock In {\em Proc. 42nd FOCS}, pp. 600--609. IEEE Comp. Soc. Press, 2001.
\newblock [\epfmtdoi{10.1109/SFCS.2001.959936}]

\bibitem{KKMO}\bibhead{KKMO}
{\sc Subhash Khot, Guy Kindler, Elchanan Mossel, and Ryan O'Donnell}: Optimal
  inapproximability results for {MAX-CUT} and other 2-variable {CSPs}?
\newblock {\em SIAM J. Comput.}, 37(1):319--357, 2007.
\newblock Preliminary version in
  \href{http://dx.doi.org/10.1109/FOCS.2004.49}{FOCS'04}. See also at
  \href{http://eccc.hpi-web.de/report/2005/101/}{ECCC}.
\newblock [\epfmtdoi{10.1137/S0097539705447372}]

\bibitem{KhotPonnuswami}\bibhead{KhotPonnuswami}
{\sc Subhash Khot and Ashok~Kumar Ponnuswami}: Better inapproximability results
  for maxclique, chromatic number and {M}in-3{L}in-{D}eletion.
\newblock In {\em Proc. 33th Internat. Colloq. on Automata, Languages and
  Programming (ICALP'06)}, pp. 226--237. Springer, 2006.
\newblock [\epfmtdoi{10.1007/11786986\_21}]

\bibitem{khotsafra}\bibhead{khotsafra}
{\sc Subhash Khot and Muli Safra}: A two prover one round game with strong
  soundness.
\newblock In {\em Proc. 52nd FOCS}, pp. 648--657. IEEE Comp. Soc. Press, 2011.
\newblock [\epfmtdoi{10.1109/FOCS.2011.62}]

\bibitem{Meg}\bibhead{Meg}
{\sc Alexandre Megretski}: Relaxation of quadratic programs in operator theory
  and system analysis.
\newblock In {\em Proc. International Workshop on Operator Theory and
  Applications (IWOTA'00)}, pp. 365--392. Springer, 2000.
\newblock [\epfmtdoi{10.1007/978-3-0348-8362-7\_15}]

\bibitem{MoshkovitzPGC}\bibhead{MoshkovitzPGC}
{\sc Dana Moshkovitz}: The projection games conjecture and the
  $\mathsf{NP}$-hardness of $\ln n$-approximating set-cover.
\newblock In {\em Proc. 15th Internat. Workshop on Approximation Algorithms for
  Combinatorial Optimization Problems (APPROX'12)}, pp. 276--287. Springer,
  2012.
\newblock See also at \href{http://eccc.hpi-web.de/report/2011/112/}{ECCC}.
\newblock [\epfmtdoi{10.1007/978-3-642-32512-0\_24}]

\bibitem{MoshRaz}\bibhead{MoshRaz}
{\sc Dana Moshkovitz and Ran Raz}: Two-query {PCP} with subconstant error.
\newblock {\em J. ACM}, 57(5):29:1--29:29, 2010.
\newblock Preliminary version in
  \href{http://dx.doi.org/10.1109/FOCS.2008.60}{FOCS'08}. See also at
  \href{http://eccc.hpi-web.de/report/2008/071/}{ECCC}.
\newblock [\epfmtdoi{10.1145/1754399.1754402}]

\bibitem{NRT}\bibhead{NRT}
{\sc Arkadi Nemirovski, Kees Roos, and Tam{\'a}s Terlaky}: On maximization of
  quadratic form over intersection of ellipsoids with common center.
\newblock {\em Mathematical Programming}, 86(3):463--473, 1999.
\newblock [\epfmtdoi{10.1007/s101070050100}]

\bibitem{Pol}\bibhead{Pol}
{\sc David Pollard}: {\em A user's guide to measure theoretic probability}.
\newblock Cambridge Univ. Press, 2002.
\newblock
  \href{http://www.cambridge.org/ar/academic/subjects/statistics-probability/probability-theory-and-stochastic-processes/users-guide-measure-theoretic-probability}{CUP}.

\bibitem{ARao}\bibhead{ARao}
{\sc Anup Rao}: Parallel repetition in projection games and a concentration
  bound.
\newblock {\em SIAM J. Comput.}, 40(6):1871--1891, 2011.
\newblock Preliminary version in
  \href{http://dx.doi.org/10.1145/1374376.1374378}{STOC'08}.
\newblock [\epfmtdoi{10.1137/080734042}]

\bibitem{Raz}\bibhead{Raz}
{\sc Ran Raz}: A parallel repetition theorem.
\newblock {\em SIAM J. Comput.}, 27(3):763--803, 1998.
\newblock Preliminary version in
  \href{http://dx.doi.org/10.1145/225058.225181}{STOC'95}.
\newblock [\epfmtdoi{10.1137/S0097539795280895}]

\bibitem{RazSafra}\bibhead{RazSafra}
{\sc Ran Raz and Shmuel Safra}: A sub-constant error-probability low-degree
  test, and a sub-constant error-probability {PCP} characterization of
  $\mathsf{NP}$.
\newblock In {\em Proc. 29th STOC}, pp. 475--484. ACM Press, 1997.
\newblock [\epfmtdoi{10.1145/258533.258641}]

\bibitem{SaT}\bibhead{SaT}
{\sc Alex Samorodnitsky and Luca Trevisan}: A {PCP} characterization of
  $\mathsf{NP}$ with optimal amortized query complexity.
\newblock In {\em Proc. 32nd STOC}, pp. 191--199. ACM Press, 2000.
\newblock [\epfmtdoi{10.1145/335305.335329}]

\end{thebibliography}

